\subsubsection{Kebutuhan Fungsional}

Sistem audit trail merupakan komponen tambahan pada arsitektur DecMed yang bertugas mengumpulkan aktivitas yang terjadi pada komponen \textit{Client}, \textit{Proxy}, maupun \textit{Storage} (\textit{on-chain} dan \textit{off-chain}), kemudian menyimpannya sebagai audit record yang bersifat \textit{append-only} dan dapat diverifikasi integritasnya. Audit record yang tersimpan selanjutnya dapat diakses melalui antarmuka pencarian oleh pihak yang berwenang untuk keperluan audit maupun investigasi insiden keamanan. Ilustrasi sederhana mengenai alur kerja sistem audit trail terhadap arsitektur DecMed digambarkan pada Gambar~\ref{fig:ilustrasi-audit-trail}.

% \begin{figure}[h]
% \centering
% \begin{tikzpicture}[
%     box/.style={draw, rectangle, minimum width=3.2cm, minimum height=1cm, align=center},
%     node distance=1.3cm and 1.6cm
% ]
%     \node[box, fill=orange!20] (client) {Client};
%     \node[box, fill=blue!15, right=of client] (proxy) {Proxy};
%     \node[box, fill=green!15, below=of client] (storage) {Storage\\(\textit{On-chain}/\textit{Off-chain})};
%     \node[box, fill=yellow!20, below=1.6cm of $(client)!0.5!(proxy)$] (collector) {Audit Trail Collector};
%     \node[box, fill=yellow!30, below=of collector] (store) {Audit Trail Store\\(\textit{append-only})};
%     \node[box, fill=purple!15, below=of store] (interface) {Antarmuka Audit Trail};
%     \node[box, fill=gray!15, below=of interface] (user) {Auditor / Pengguna};

%     \draw[->] (client) -- node[above]{HTTP request} (proxy);
%     \draw[->] (client) |- (collector);
%     \draw[->] (proxy) |- (collector);
%     \draw[->] (storage) |- (collector);
%     \draw[->] (client) -- (storage);
%     \draw[->] (proxy) -- (storage);
%     \draw[->] (collector) -- node[right]{Simpan \textit{event}} (store);
%     \draw[->] (store) -- node[right]{Kueri} (interface);
%     \draw[->] (interface) -- node[right]{Akses hasil} (user);
% \end{tikzpicture}
% \caption{Ilustrasi Alur Kerja Sistem Audit Trail}
% \label{fig:ilustrasi-audit-trail}
% \end{figure}

Berdasarkan alur kerja tersebut, ditentukan aktor yang terlibat pada sistem audit trail beserta \textit{privileges} yang dimiliki masing-masing aktor, sebagaimana disajikan pada Tabel~\ref{tab:aktor-privileges-audit-trail}.

\begin{xltabular}{\textwidth}{|>{\raggedright\arraybackslash}p{3.2cm}|X|}
    \caption{Aktor dan Privileges Sistem Audit Trail} \\
    \hline
    \textbf{Aktor} & \textbf{Privileges} \\
    \hline
    \endfirsthead
    \hline
    \textbf{Aktor} & \textbf{Privileges} \\
    \hline
    \endhead
    \hline
    \endfoot
    \hline
    \endlastfoot

    Sistem &
    \begin{itemize}
        \item Membangkitkan audit record secara otomatis untuk setiap aktivitas yang terjadi pada komponen DecMed,
        \item Menyimpan audit record secara \textit{append-only}, dan
        \item Menerapkan kebijakan retensi terhadap audit record yang tersimpan.
    \end{itemize} \\
    \hline

    Auditor &
    \begin{itemize}
        \item Mencari dan menyaring audit record berdasarkan jenis aktivitas, aktor, rentang waktu, dan objek terkait,
        \item Melihat detail audit record,
        \item Memverifikasi integritas rangkaian audit record, dan
        \item Mengekspor audit record untuk keperluan investigasi.
    \end{itemize} \\
    \hline

    Admin Sistem &
    \begin{itemize}
        \item Mengatur kebijakan retensi audit record, dan
        \item Mengatur hak akses pengguna terhadap audit trail.
    \end{itemize} \\
    \hline

    Pasien &
    \begin{itemize}
        \item Melihat audit record terkait aktivitas pada data RME miliknya.
    \end{itemize} \\
    \hline

    Tenaga Kesehatan &
    \begin{itemize}
        \item Melihat audit record terkait aktivitas yang ia lakukan terhadap data RME pasien.
    \end{itemize} \\
    \hline

    Petugas Rekam Medis &
    \begin{itemize}
        \item Melihat audit record terkait aktivitas yang ia lakukan terhadap data administratif pasien.
    \end{itemize} \\
    \hline

    Admin Fasyankes &
    \begin{itemize}
        \item Melihat audit record terkait aktivitas penerbitan kunci aktivasi yang ia lakukan.
    \end{itemize} \\
    \hline

    Kementerian Kesehatan &
    \begin{itemize}
        \item Melihat audit record terkait aktivitas registrasi fasyankes yang ia lakukan.
    \end{itemize} \\
    \hline

\end{xltabular}

Berdasarkan aktor dan \textit{privileges} tersebut, disusun daftar kebutuhan fungsional sistem audit trail sebagaimana disajikan pada Tabel~\ref{tab:kebutuhan-fungsional}.

\begin{xltabular}{\textwidth}{|>{\centering\arraybackslash}p{1.3cm}|X|}
    \caption{Kebutuhan Fungsional Sistem Audit Trail} \\
    \hline
    \textbf{ID} & \textbf{Deskripsi} \\
    \hline
    \endfirsthead
    \hline
    \textbf{ID} & \textbf{Deskripsi} \\
    \hline
    \endhead
    \hline
    \endfoot
    \hline
    \endlastfoot

    F01 & Sistem dapat membangkitkan audit record secara otomatis setiap kali \textit{event} pada komponen DecMed terjadi \\
    \hline
    F02 & Sistem dapat menyimpan audit record secara \textit{append-only} sehingga audit record yang telah dibuat tidak dapat diubah maupun dihapus \\
    \hline
    F03 & Sistem dapat menyusun audit record ke dalam rangkaian \textit{hash} untuk menjaga integritas antar audit record \\
    \hline
    F04 & Sistem dapat mencatat audit record baik untuk aktivitas yang berhasil maupun gagal/ditolak \\
    \hline
    F05 & Sistem dapat menerapkan kebijakan retensi audit record sesuai jangka waktu yang ditentukan \\
    \hline
    F06 & Auditor dapat melakukan pencarian audit record berdasarkan jenis aktivitas \\
    \hline
    F07 & Auditor dapat melakukan pencarian audit record berdasarkan identitas aktor \\
    \hline
    F08 & Auditor dapat melakukan pencarian audit record berdasarkan rentang waktu \\
    \hline
    F09 & Auditor dapat melakukan pencarian audit record berdasarkan objek yang terkait \\
    \hline
    F10 & Auditor dapat melihat detail lengkap suatu audit record \\
    \hline
    F11 & Auditor dapat memverifikasi integritas rangkaian audit record \\
    \hline
    F12 & Auditor dapat mengekspor audit record ke dalam format standar \\
    \hline
    F13 & Admin sistem dapat mengatur kebijakan retensi audit record \\
    \hline
    F14 & Admin sistem dapat mengatur hak akses pengguna terhadap audit trail \\
    \hline
    F15 & Pasien dapat melihat audit record terkait aktivitas pada data RME miliknya \\
    \hline
    F16 & Tenaga kesehatan dapat melihat audit record terkait aktivitas yang ia lakukan terhadap data RME pasien \\
    \hline
    F17 & Petugas rekam medis dapat melihat audit record terkait aktivitas yang ia lakukan terhadap data administratif pasien \\
    \hline
    F18 & Admin fasyankes dapat melihat audit record terkait aktivitas penerbitan kunci aktivasi yang ia lakukan \\
    \hline
    F19 & Kementerian kesehatan dapat melihat audit record terkait aktivitas registrasi fasyankes yang ia lakukan \\
    \hline

\end{xltabular}